# Chinchilla-Lutar Scaling Laws with Grokking Phase Detection, E8-Triality Mixture-of-Experts, and Dynamical Systems Bifurcation Analysis

**Author:** Stephen Paul Lutar Jr.
**Affiliation:** SZL Consulting Ltd
**ORCID:** [0009-0001-0110-4173](https://orcid.org/0009-0001-0110-4173)
**Contact:** stephenlutar2@gmail.com
**Date:** 4 May 2026
**Version:** v11
**Companion papers:** v1--v10 (see References)
**Reference implementation:** [github.com/szl-holdings/szl-holdings-platform](https://github.com/szl-holdings/szl-holdings-platform) -- `packages/ouroboros-integrations/src/sovereign-engine.ts` (CLS, GPD, LSR, E8, HQO, CMST, DSBD classes)
**License:** CC BY 4.0

---

## Abstract

We present four interconnected contributions to AI training dynamics and model architecture: (1) Chinchilla-Lutar Scaling (CLS), which extends the Chinchilla scaling law (Hoffmann et al., 2022) with Omega-weighted compute allocation; (2) Grokking Phase Detection (GPD), which detects memorization-to-generalization phase transitions in real time; (3) E8-Triality Mixture-of-Experts (E8-MoE) with the Lutar Simplex Router, which routes queries to expert slots using the 192-slot E8 lattice; and (4) Dynamical Systems Bifurcation Detection (DSBD), which classifies training dynamics bifurcations (saddle-node, Hopf, period-doubling, transcritical) and distinguishes integrator from resonator neural regimes. We also cover the Condor Mamba-SSM State Tracker (CMST) for linear-time state tracking and the Hilbert QAOA-Omega (HQO) variational optimizer. This final paper in the thesis lineage (v11) consolidates the architectural and training dynamics innovations.

This paper does **not** claim empirical validation against standard scaling law benchmarks. The contribution is the formal framework and its public reference implementation.

---

## 1. Motivation and Continuity from v3--v10

The v3--v10 papers established the trust, quality, retrieval, safety, memory, inference, interpretability, and quantum-coherence layers. The present paper addresses the *training and architecture* substrate: how to allocate compute, detect phase transitions, route experts, and analyze dynamical stability.

---

## 2. Chinchilla-Lutar Scaling (CLS)

### 2.1 Standard Chinchilla

Hoffmann et al. (2022) established that for a compute budget \( C \) FLOPs:

\[
N^* \propto C^{0.5}, \quad D^* \propto C^{0.5}
\]

where \( N^* \) is optimal parameter count and \( D^* \) is optimal token count.

### 2.2 Omega-Weighted Extension

CLS extends this with the Lutar Omega correction:

\[
N^*_\Omega = N^* \cdot L_\Omega^{1/4}, \quad D^*_\Omega = D^* \cdot L_\Omega^{1/4}
\]

When \( L_\Omega = 1 \), standard Chinchilla is recovered. When \( L_\Omega > 1 \) (high-quality evaluation context), more parameters and data are allocated. When \( L_\Omega < 1 \), compute is conserved.

### 2.3 Inference-Time Adjustment

For an inference budget \( I \) tokens/second:

\[
\text{batch}_\Omega = \lfloor I \cdot (1 - L_\Omega \cdot H) \rfloor
\]

where \( H \) is the Hermetic temperature (v6), trading inference throughput against quality.

---

## 3. Grokking Phase Detection (GPD)

### 3.1 Background

Grokking (Power et al., 2022) refers to the delayed generalization phenomenon where a model memorizes training data, then suddenly transitions to genuine generalization after continued training.

### 3.2 Phase Classification

GPD monitors three signals: gradient norm, weight entropy, and synergy (gradient-entropy coupling):

\[
\text{synergy} = \frac{|\nabla \theta| \cdot H(\theta)}{|\nabla \theta| + H(\theta)}
\]

Four phases are classified:

| Phase | Gradient Norm | Weight Entropy | Synergy | Interpretation |
|---|---|---|---|---|
| Memorization | High | Low | Low | Fitting training data |
| Circuit formation | Declining | Rising | Medium | Building generalizing structure |
| Grokking | Low | High | High | Phase transition to generalization |
| Stable generalization | Low | Stable | Stable | Post-transition equilibrium |

### 3.3 Transition Prediction

GPD predicts the grokking transition by extrapolating the synergy trajectory:

\[
\hat{t}_{\text{grok}} = t + \frac{s_{\text{threshold}} - s_t}{\Delta s / \Delta t}
\]

---

## 4. E8-Triality Mixture-of-Experts

### 4.1 E8 Lattice Structure

The \( E_8 \) root system has 240 roots in 8 dimensions. After removing the 48 Weyl-orbit degenerate roots, 192 slots remain. The triality automorphism cycles through three 8-dimensional representations.

### 4.2 Lutar Simplex Router (LSR)

The LSR routes a query \( q \) to an expert slot:

1. Compute query complexity: \( c(q) = |\text{tokens}(q)| \cdot \text{entropy}(q) \)
2. Map to E8 slot: \( \text{slot} = \text{hash}(q) \bmod 192 \)
3. Select triality representation: \( \text{rep} = \lfloor \text{slot} / 64 \rfloor \bmod 3 \)
4. Select provider by complexity threshold

### 4.3 Hilbert QAOA-Omega (HQO)

The HQO extends the Quantum Approximate Optimization Algorithm with Omega weighting:

\[
|\psi(\gamma, \beta)\rangle = \prod_{p=1}^{P} e^{-i\beta_p B} e^{-i\gamma_p C} |+\rangle^{\otimes n}
\]

The objective function is the Omega-weighted expected value:

\[
\langle C \rangle_\Omega = L_\Omega \cdot \langle \psi(\gamma, \beta) | C | \psi(\gamma, \beta) \rangle
\]

---

## 5. Condor Mamba-SSM State Tracker (CMST)

### 5.1 State-Space Model

The CMST implements a complex-valued linear state-space model:

\[
\mathbf{h}_{t+1} = A \mathbf{h}_t + B x_t, \quad y_t = C \mathbf{h}_t + D x_t
\]

where \( A \in \mathbb{C}^{d \times d} \) is the state transition matrix initialized with HiPPO-LegS (Gu et al., 2020) or S4 diagonal parameterization (Gu et al., 2021).

### 5.2 Selective State-Space (Mamba)

Following Gu and Dao (2023), CMST adds input-dependent selection:

\[
A_t = A + \Delta A(x_t), \quad B_t = B + \Delta B(x_t)
\]

enabling the model to selectively propagate or forget state based on input content.

---

## 6. Dynamical Systems Bifurcation Detector (DSBD)

### 6.1 Background

Dynamical systems theory (Strogatz, 2015) classifies bifurcations --- qualitative changes in system behavior as a parameter varies. We apply this framework to training dynamics.

### 6.2 Bifurcation Types

| Type | Eigenvalue Signature | Neural Interpretation |
|---|---|---|
| Saddle-node | Real eigenvalue crosses zero | Loss landscape critical point creation/annihilation |
| Hopf | Complex eigenvalue pair crosses imaginary axis | Oscillatory training dynamics onset |
| Period-doubling | Real eigenvalue crosses -1 | Training oscillation period doubling |
| Transcritical | Two real eigenvalues exchange stability | Regime switching between training modes |

### 6.3 Integrator vs. Resonator Classification

Following Izhikevich (2007) and Kirsanov (2024):

- **Integrator neurons:** Respond to sustained input. Real eigenvalues. Saddle-node bifurcation.
- **Resonator neurons:** Respond to oscillatory input. Complex eigenvalues. Hopf bifurcation.

Classification rule: if \( |\text{Im}(\lambda)| > 0.1 \), the regime is resonator; otherwise integrator.

### 6.4 Lyapunov Exponent Estimation

\[
\lambda_{\text{Lyap}} = \lim_{t \to \infty} \frac{1}{t} \ln \frac{|\delta \mathbf{x}(t)|}{|\delta \mathbf{x}(0)|}
\]

Positive Lyapunov exponents indicate chaotic training dynamics. Negative exponents indicate convergent training.

### 6.5 Upcoming Bifurcation Detection

The DSBD predicts upcoming bifurcations by extrapolating the eigenvalue trajectory:

\[
\hat{t}_{\text{bif}} = t + \frac{-\text{Re}(\lambda_t)}{\Delta \text{Re}(\lambda) / \Delta t}
\]

If \( \hat{t}_{\text{bif}} \) is within a lookahead window, a warning is issued.

---

## 7. Integration with the Full Thesis Lineage

| Version | Focus | Training/Architecture Role |
|---|---|---|
| v3 | Trust aggregation | Quality metric for training checkpoints |
| v4 | Omega Formalism | Compute allocation via CLS |
| v5 | Retrieval | Training data selection via Prisca-GraphRAG |
| v6 | Safety | Guard on training data via HCG |
| v7 | Memory | Catastrophic forgetting prevention via SCL |
| v8 | Active inference | Training as free-energy minimization via FELAI |
| v9 | Interpretability | Feature monitoring via TSA |
| v10 | Quantum/geometry | Correlation structure via EBEV/SGCE |
| v11 (this) | Training dynamics | CLS, GPD, E8-MoE, CMST, DSBD |

---

## 8. What this paper does and does not claim

**Claims:**
1. CLS is a well-defined extension of the Chinchilla scaling law with Omega weighting.
2. GPD phase classification follows from the synergy metric definition.
3. The E8 routing table uses the 192-slot reduced root system.
4. DSBD bifurcation classification follows standard dynamical systems theory (Strogatz, 2015).
5. All components are implemented in the public reference implementation.

**Non-claims:**
- No claim of empirical superiority over standard scaling laws on LLM training runs.
- No claim that grokking prediction is accurate on standard modular arithmetic benchmarks.
- No claim that E8 routing improves MoE performance over standard hash-based routing.
- No claim of deployed production use.

---

## 9. Conclusion

This final paper in the Ouroboros Thesis lineage (v11) consolidates the training dynamics and architectural innovations: Chinchilla-Lutar scaling, grokking phase detection, E8-triality mixture-of-experts, Mamba state-space tracking, and dynamical systems bifurcation analysis. Together with v3--v10, the eleven-paper series provides a comprehensive formal framework spanning trust, quality, retrieval, safety, memory, inference, interpretability, quantum foundations, and training dynamics.

---

## References

1. Hoffmann, J. et al. (2022). Training compute-optimal large language models. *NeurIPS*.
2. Power, A. et al. (2022). Grokking: Generalization beyond overfitting on small algorithmic datasets. *ICLR*.
3. Strogatz, S.H. (2015). *Nonlinear Dynamics and Chaos*. Westview Press, 2nd edition.
4. Izhikevich, E.M. (2007). *Dynamical Systems in Neuroscience*. MIT Press.
5. Gu, A., Goel, K., Re, C. (2021). Efficiently modeling long sequences with structured state spaces. *ICLR 2022*.
6. Gu, A. and Dao, T. (2023). Mamba: Linear-time sequence modeling with selective state spaces. *arXiv:2312.00752*.
7. Kirsanov, A. (2024). Dynamical systems in neuroscience trilogy. YouTube lecture series.
8. Fedus, W. et al. (2022). Switch Transformers: Scaling to trillion parameter models with simple and efficient sparsity. *JMLR*.
9. Lutar, S.P. (2026a--j). Ouroboros Thesis series v1--v10. See companion papers.
